% ---------------------------------------------------------------------------
% ccscreen.tex -- character-cell screen mock-ups, drawn in TikZ
%
% Everything the manual shows of the display is built out of character cells
% on the same 80-column geometry the program uses, in the Commander X16's own
% palette. A cell is \ccchw wide and \ccchh tall; the picture's coordinate
% system is one unit per cell with the origin at the top left, so a figure can
% be written with the same row and column numbers the code uses.
%
% Spaces inside text fragments must be written ~ , which is exactly one cell
% wide in a monospaced font and does not collapse the way a run of ordinary
% spaces would.
% ---------------------------------------------------------------------------

\newlength{\ccchw}      % width of one character cell
\newlength{\ccchh}      % height of one character cell
\newlength{\ccscrfs}    % font size used inside the cells

% Scale factor 1 gives a full 80-column screen just inside the type block.
\newcommand{\ccscreenscale}[1]{%
  \setlength{\ccchw}{#1\dimexpr 4.2pt\relax}%
  \setlength{\ccchh}{#1\dimexpr 5.4pt\relax}%
  \setlength{\ccscrfs}{#1\dimexpr 7.0pt\relax}%
}
\ccscreenscale{1}

% The X16 attribute byte is (background << 4) | foreground. These are the
% combinations the program actually uses, named as the source names them.
\newcommand{\ccattr}[1]{\csname ccattr@#1\endcsname}
\expandafter\def\csname ccattr@menu@bg\endcsname{x16lgrey}
\expandafter\def\csname ccattr@menu@fg\endcsname{x16black}

\newcommand{\ccscreenbegin}[2]{%
  \begin{tikzpicture}[x=\ccchw,y=-\ccchh,
      cctxt/.style={anchor=base west,inner sep=0pt,outer sep=0pt,
                    font=\ttfamily\fontsize{\the\ccscrfs}{\the\ccscrfs}\selectfont}]
    \def\ccscrcols{#1}\def\ccscrrows{#2}%
    % Clip to the screen. The real display has the same edge: a column that
    % starts before column 80 and runs past it is cut off there rather than
    % being left out, and text is never allowed to leave its cell.
    \clip (0,0) rectangle (#1,#2);
    % the paper: every cell starts white, as screen_clear leaves it
    \fill[x16white] (0,0) rectangle (#1,#2);
}
\newcommand{\ccscreenend}{%
    \draw[ccink,line width=0.7pt] (0,0) rectangle (\ccscrcols,\ccscrrows);
  \end{tikzpicture}%
}

% --- fills -----------------------------------------------------------------
% \ccfill{row}{col}{width}{colour}
\newcommand{\ccfill}[4]{\fill[#4] (#2,#1) rectangle ++(#3,1);}
% \ccfillrow{row}{colour}
\newcommand{\ccfillrow}[2]{\fill[#2] (0,#1) rectangle ++(\ccscrcols,1);}
% \ccfillbox{row}{col}{width}{height}{colour}
\newcommand{\ccfillbox}[5]{\fill[#5] (#2,#1) rectangle ++(#3,#4);}

% --- text ------------------------------------------------------------------
% The baseline sits 0.78 of a cell down, which is where a character sits in an
% 8-pixel tile.
% \cctext{row}{col}{colour}{text}
\newcommand{\cctext}[4]{%
  \node[cctxt,text=#3,yshift=0.22\ccchh] at (#2,\the\numexpr#1+1\relax) {#4};}

% \cctextr{row}{col}{width}{colour}{text} -- right-aligned in col..col+width-1
\newcommand{\cctextr}[5]{%
  \node[cctxt,anchor=base east,text=#4,yshift=0.22\ccchh]
    at (\the\numexpr#2+#3\relax,\the\numexpr#1+1\relax) {#5};}

% --- worksheet geometry ----------------------------------------------------
% GRID_ROWHDR_W is 6: five digits of row number and one separator column.
% Columns are X16S_DEF_COL_W = 10 cells wide unless a file says otherwise.
\newcommand{\ccgridx}[1]{\the\numexpr 6+10*(#1)\relax}

% \ccrowhdr{row}{label}{fg}{bg} -- the row-number gutter
\newcommand{\ccrowhdr}[4]{%
  \ccfill{#1}{0}{6}{#4}%
  \cctextr{#1}{0}{5}{#3}{#2}}

% \cccellL{row}{colindex}{fg}{text} -- a left-aligned cell value
\newcommand{\cccellL}[4]{\cctext{#1}{\ccgridx{#2}}{#3}{~#4}}
% \cccellR{row}{colindex}{fg}{text} -- a right-aligned cell value
\newcommand{\cccellR}[4]{\cctextr{#1}{\ccgridx{#2}}{10}{#3}{#4}}

% \cccursor{row}{colindex} -- the cell cursor, C_CURSOR is reverse video
\newcommand{\cccursor}[2]{\ccfill{#1}{\ccgridx{#2}}{10}{x16black}}

% \cccolhdr{row}{colindex}{letter}{selected?}
% The branch is taken outside \ccfill: a conditional left inside a TikZ option
% list is expanded while the key parser is picking the argument apart, and
% that goes wrong in ways that are not worth debugging twice.
\newcommand{\cccolhdr}[4]{%
  \ifnum#4=1\relax
    \ccfill{#1}{\ccgridx{#2}}{10}{x16yellow}%
  \else
    \ccfill{#1}{\ccgridx{#2}}{10}{x16lgrey}%
  \fi
  \cctext{#1}{\ccgridx{#2}}{x16black}{~#3}}

% --- callouts --------------------------------------------------------------
% A label outside the artwork with a hairline leader to a point on it, the way
% every screen figure in a manual of this period is annotated.
% \cccallout[anchor]{label}{label coord}{target coord}
% Both coordinates are written with their own parentheses, so that a named
% node or a calc expression can be used as freely as a pair of numbers.
\newcommand{\cccallout}[4][west]{%
  \draw[ccink,line width=0.4pt] #3 -- #4;
  \fill[ccink] #4 circle (0.4);
  \node[anchor=#1,inner sep=2.5pt,
        font=\sffamily\fontsize{7.4pt}{8.6pt}\selectfont,text=ccink,align=left] at #3 {#2};}

% --- dialogue boxes --------------------------------------------------------
% DLG_X is 12 and DLG_W is 52, so every dialogue in the program occupies
% columns 12 to 63. The title row is C_TITLE, the body C_DLG.
\newcommand{\ccdlgbox}[3]{%
  \ccfill{#1}{12}{52}{x16blue}%
  \cctext{#1}{14}{x16white}{#3}%
  \ccfillbox{\the\numexpr#1+1\relax}{12}{52}{\the\numexpr#2-1\relax}{x16lgrey}}
